% ============================================================
% PASTE THIS INSIDE \subsection{Accumulator Variables}
% Replace the single placeholder line:
%   "Describe the accumulator role, the evidence supporting..."
% ============================================================

\label{sec:accum}
For accumulator variables, we probe Qwen2.5-1.5B residual activations on XLCoST
programs: 500 Python and 300 each of Java, C++, C, and C\#. We identify
accumulator variables via Python's \texttt{ast} module and language-specific
regex patterns, detecting variables targeted by \texttt{+=}/\texttt{-=} or
\texttt{.append()}/\texttt{.extend()} inside loop bodies. Hidden states are
mean-pooled over each variable's token span and fed to a layer-wise logistic
regression probe (class-balanced, $C{=}1.0$, 80/20 split).

To separate lexical from structural signal, we apply six variable-naming
strategies to the same 500 Python programs; results are in
Table~\ref{tab:accum-perturbation} and
Figures~\ref{fig:appendix-sree-accum-probe}--\ref{fig:appendix-sree-accum-delta}.
The baseline achieves F1~$= 0.893$ at L9, lower than the index role ($0.959$)
because accumulator names (\texttt{total}, \texttt{count}, \texttt{result})
overlap heavily with non-accumulator identifiers. Counterintuitively, removing
or adversarially swapping names \emph{improves} accuracy: all-same naming
($+0.076$) and misleading index-like names ($+0.079$) both exceed the baseline,
as the structural cue---being the \texttt{+=} target inside a loop---becomes the
only signal left. This is the \emph{opposite} of the index role. The numeric
strategy peaks at L3, far earlier than other conditions (L5--L10), suggesting
opaque naming pushes role encoding into shallower layers.

\begin{table}[t]
\centering
\caption{Accumulator probe results under six naming strategies, Python XLCoST,
Qwen2.5-1.5B.}
\label{tab:accum-perturbation}
\begin{tabular}{lcccc}
\toprule
Strategy & Best Layer & Test Acc & Test F1 & $\Delta$ F1 \\
\midrule
Baseline (original names)           & L9  & 0.984 & 0.893 & ---     \\
Random nouns                        & L10 & 0.979 & 0.870 & $-$0.023 \\
Single chars (\texttt{a, b, c\ldots}) & L9  & 0.971 & 0.840 & $-$0.053 \\
All same (\texttt{x})               & L5  & 0.986 & 0.969 & $+$0.076 \\
Numeric (\texttt{v1, v2\ldots})     & L3  & 0.999 & 0.986 & $+$0.093 \\
Misleading (index-style names)      & L9  & 0.996 & 0.972 & $+$0.079 \\
\bottomrule
\end{tabular}
\end{table}

Cross-language transfer confirms the representation is language-agnostic
(Table~\ref{tab:accum-crosslang}, Figure~\ref{fig:appendix-sree-accum-crosslang}).
A Python-trained probe (L9) achieves $0.943$--$0.948$ accuracy on Java, C++, C,
and C\# without seeing any compiled-language programs. In-domain F1 for those
languages ($0.939$--$0.970$) consistently exceeds Python ($0.893$), likely
because \texttt{+=} in a compiled-language loop body is less syntactically
ambiguous than Python's mix of \texttt{+=} and \texttt{.append()} idioms.

\begin{table}[t]
\centering
\caption{Cross-language accumulator probe. Python-trained probe (L9) applied
to four additional languages.}
\label{tab:accum-crosslang}
\begin{tabular}{lcccc}
\toprule
Language & Programs & In-Domain F1 & Cross-Lang Acc \\
\midrule
Python     & 181 / 500 & 0.893 (L9)  & ---   \\
Java       & 208 / 300 & 0.952 (L9)  & 0.948 \\
C++        & 204 / 300 & 0.939 (L7)  & 0.948 \\
C          & 165 / 300 & 0.970 (L11) & 0.943 \\
C\#        & 207 / 300 & 0.950 (L9)  & 0.945 \\
\bottomrule
\end{tabular}
\end{table}
